\section{Discussion}
\label{sec:discussion}

\subsection{What survives of the original claim}

The premise that motivated this work was attractive: an irrigation command is
already a digital event, so MRV evidence should be nearly free. That premise
survives, but not in the form we first published. It survives because the
controller that issues the command can also sign it, and the signed record is
what makes the evidence admissible --- not because a hydrological model can
detect a lie. Where the original draft claimed that physics attests the log, the
correct statement is that physics \emph{constrains} the log: it rules out depths
no soil can hold and states no field can reach, which is worth having, and it is
silent about the one transformation that pays.

\subsection{Why publishing the threshold is not the problem}

One might conclude that $\varepsilon$ should be kept secret. It should not, for
two reasons. A verifier must be able to reproduce a decision, which requires
published parameters. And secrecy would not help: by Proposition~\ref{prop:blind}
the time-shifting attack preserves the total volume and every individual depth,
so it does not need to know $\varepsilon$ in order to stay beneath the
infiltration bound --- it needs only to avoid moving water, which is visible to
any farmer. The vulnerability is structural, not a consequence of disclosure.

\subsection{Relation to remote sensing}

Satellite radar detects irrigation events but revisits every two to four days and
cannot resolve a $72$\,h dry phase in mid-period, which is the crediting unit.
It therefore cannot substitute for a command record, and it cannot detect
time-shifting within a revisit interval either. Its role remains the one VM0051
Appendix 4 assigns it: a per-plot cross-check that lowers the frequency of field
verification~\cite{vm0051}.

\subsection{Limitations}

\begin{enumerate}
\item \emph{Simulated farmers.} Forgeries are injected into simulated reference
trajectories. A field pilot is needed to measure behavioural response, in
particular whether farmers shift their \emph{actual} pumping once the anchor is
known to exist --- an attack our model cannot represent, since it changes $M$
rather than $L$.
\item \emph{The anchor assumes an honest gateway.}
Proposition~\ref{prop:complete} is completeness against log forgery, not against
a compromised controller. Device integrity, key custody and replay protection
are assumed from the companion work and are not re-established here.
\item \emph{Not applicable to manual irrigation.} Where a sluice is opened by
hand there is no command record, so the anchor is unavailable. This excludes a
substantial share of delta farms and bounds the addressable population of the
scheme.
\item \emph{$\Theta$ is a judgement.} Five textures, two $K_c$ regimes, three
efficiencies and five initial states bracket the delta but do not exhaust it. A
field outside $\Theta$ loses the guarantee of Proposition~\ref{prop:noalarm},
which is exactly the situation the proportional threshold and the query verdict
exist to absorb.
\item \emph{Device prices are planning estimates, not quotations.} The tier
ordering and the break-even inequality are unaffected; the absolute margin is.
\item \emph{No emission quantification.} The layer produces activity data.
Coupling it to a methane model and locally calibrated emission factors is
separate work, as is the cost table's assumption that credits scale with
$n_{\mathrm{dry}}$.
\end{enumerate}

\section{Conclusion}

We set out to make rice MRV nearly free by auditing a farmer-written irrigation
log against soil-water physics, and the first version of that layer reported no
false accusation and complete detection. Both numbers were artefacts of a field
generated by the auditing model itself. Rebuilding the field along seven axes of
realism, and replacing fixed forgery recipes with an adversary that optimises the
credited quantity, changes the conclusion.

The point-model detector falsely accuses $26.7\%$ of honest farmers, because an
absolute pump-gap threshold cannot distinguish a forgotten entry from a secret
irrigation. A proportional threshold and an intersection over a $150$-element
uncertainty set reduce that to $5.0\%$ and $0/60$ respectively, at a measured
cost in detection of omission fraud, from $78.3\%$ to $40\%$. Meanwhile a
volume-preserving, depth-bounded permutation of claimed times inflates creditable
dry phases by $69\%$ to $78\%$ while escaping both physical tiers in $91\%$ to
$100\%$ of cases, and we prove this is unavoidable for any detector whose verdict
is a function of the replayed trajectory: delaying a claim both lengthens the dry
phase and increases the deficit available to it, so profit and undetectability
share a gradient.

The repair is not hydrological. Matching the claimed log against the gateway's
signed, time-stamped command record catches $116/116$ genuinely forged logs and
accuses none of $60$ honest ones, and it costs nothing beyond a controller the
companion work already deploys. The honest formulation of the result is
therefore narrower and more useful than the one we started with: near-zero-cost
MRV for low-emission rice is achievable, but only where a device stands between
the farmer and the pump. Three previously published numbers are corrected
accordingly --- the admissible depth is the range $[49.0,86.0]$\,mm rather than a
single $61$\,mm; the third physical filter cannot activate on dry-season delta
rice, clipped depletion peaking at $0.9987$ over $150$ configurations; and a
subthreshold-based absolute gap test is inadmissible. Code, both uncertainty
sets, all $281$ audited logs and $46$ unit tests, one per design hole, are
public.
